\chapter{INTRODUCTION}

\section{Project Overview}
AIRES (AI-Based Resume Analyzer \& Job Match System) is a modern, full-stack web application built on the MERN stack (MongoDB, Express.js, React, Node.js) that automates and enhances the traditional recruitment workflow. The system is designed to eliminate the inefficiencies of manual resume screening by providing AI-powered resume parsing, intelligent job matching, and structured applicant tracking for both job seekers and hiring managers.\\\\
In the contemporary job market, recruiters routinely receive hundreds of applications for a single position, making it operationally challenging to review each resume manually. Simultaneously, candidates often lack visibility into how well their skill profile aligns with a given job description. AIRES addresses both of these challenges through a unified digital platform that provides data-driven insights to all stakeholders.

\section{Motivation}
The increasing volume of job applications combined with the growing complexity of skill requirements has created a significant bottleneck in the recruitment process. Traditional Applicant Tracking Systems (ATS) are often opaque, expensive, and inaccessible to smaller organizations. There is a clear need for an intelligent, open, and user-friendly system that:
\begin{itemize}
    \item Automatically extracts structured information from unstructured resume documents.
    \item Computes objective compatibility scores between a candidate's profile and a job's requirements.
    \item Provides candidates with actionable skill gap analysis.
    \item Gives recruiters a ranked, filterable view of all applicants.
    \item Enables transparent, real-time tracking of application status.
\end{itemize}
AIRES was conceived to fulfill these requirements through a clean, accessible web interface backed by AI-driven logic.

\section{Objectives}
The primary objectives of the AIRES project are:
\begin{enumerate}
    \item To develop a MERN-based web application supporting two distinct user roles: Candidate and Recruiter.
    \item To implement AI-powered resume parsing capable of extracting skills, education, and experience from uploaded PDF/DOCX files.
    \item To design a job-matching algorithm that calculates a percentage-based compatibility score between a candidate's skill set and a job's required skills.
    \item To present candidates with a ranked list of job openings ordered by their match score, along with a skill gap analysis for each position.
    \item To provide recruiters with a dashboard for posting jobs, managing applications, ranking candidates, and accessing recruitment analytics.
    \item To enable real-time application status tracking for both candidates and recruiters.
\end{enumerate}

\section{Scope}
The scope of the AIRES project encompasses the following functionalities:
\begin{itemize}
    \item \textbf{User Authentication:} Secure registration and login for both candidates and recruiters with role-based access control.
    \item \textbf{Resume Management:} Upload, parse, and store resume data; extraction of skills, education, and experience fields.
    \item \textbf{Job Management:} Recruiters can create, edit, and delete job postings with associated skill requirements.
    \item \textbf{AI Job Matching:} Automated computation of match scores using skill-based comparison algorithms.
    \item \textbf{Candidate Ranking:} Recruiter-facing ranked view of all applicants per job position.
    \item \textbf{Skill Gap Analysis:} Identification of missing skills required for a given job relative to the candidate's profile.
    \item \textbf{Application Tracking:} Candidates can apply for jobs and monitor application status; recruiters can update statuses.
    \item \textbf{Analytics Dashboard:} Recruiters have access to recruitment statistics including total jobs, applicants, shortlist rate, and pending reviews.
\end{itemize}

\section{Report Organization}
The remainder of this report is organized as follows. Chapter 2 reviews relevant literature and supporting technologies. Chapter 3 presents the system analysis including module descriptions, business rules, feasibility analysis, and system environment. Chapter 4 covers the system design with use case diagrams, activity diagrams, and detailed design decisions. Chapter 5 documents the testing strategy and test cases. Chapter 6 describes the deployment process and folder structure. Chapter 7 presents the Git history. Chapter 8 provides conclusions, followed by future work and appendices.

